% CV — Ben Lonnqvist. Compile with:  xelatex Ben_Lonnqvist_CV.tex  (run twice)
% Requires the IBM Plex fonts (free: https://github.com/IBM/plex) installed on your system.
\documentclass[10pt]{article}
\usepackage[a4paper,margin=16mm,top=14mm,bottom=15mm]{geometry}
\usepackage{fontspec}
\usepackage{xcolor}
\usepackage{needspace}
\usepackage{enumitem}
\usepackage{tabularx}\usepackage{array}
\usepackage{fancyhdr}
\usepackage{lastpage}
\usepackage[hidelinks]{hyperref}

\setmainfont{IBM Plex Sans}[UprightFont=* , BoldFont=* SemiBold, ItalicFont=* Italic, BoldItalicFont=* SemiBold Italic]
\setmonofont{IBM Plex Mono}[BoldFont=* Medium]
\newfontfamily\cond{IBM Plex Sans Condensed}[UprightFont=* SemiBold, BoldFont=* SemiBold]

% Site palette, deepened where needed for legibility on white paper
\definecolor{ink}{HTML}{101B24}     % midnight
\definecolor{muted}{HTML}{4E626C}   % darker blue-gray
\definecolor{signal}{HTML}{305D56}  % deep teal (links, markers)
\definecolor{signalmid}{HTML}{7DBFA6} % mid seafoam (second marker pixel)
\definecolor{soft}{HTML}{C0CED4}    % blue-gray rules
\color{ink}
\hypersetup{colorlinks=true,urlcolor=ink,linkcolor=ink}

\setlength{\parindent}{0pt}
\setlength{\parskip}{0pt}
\linespread{1.08}

% Section headings: magenta square + mono uppercase label + ink rule
\newcommand{\cvsection}[2][]{\par\vspace{9pt}\needspace{4\baselineskip}%
  {\ttfamily\small\textcolor{signal}{\rule{6pt}{6pt}}\textcolor{signalmid}{\rule{6pt}{6pt}}\hspace{6pt}{\addfontfeatures{LetterSpace=8}\MakeUppercase{#2}}\hfill{\rmfamily\footnotesize\color{muted}#1}}\par
  \vspace{2pt}{\color{ink}\rule{\linewidth}{1.2pt}}\par\vspace{6pt}}
\newcommand{\cvsub}[1]{\par\vspace{5pt}{\ttfamily\footnotesize\color{muted}\addfontfeatures{LetterSpace=6}\MakeUppercase{#1}}\par\vspace{3pt}}

\pagestyle{fancy}\fancyhf{}
\renewcommand{\headrulewidth}{0pt}
\fancyfoot[L]{\ttfamily\scriptsize\color{muted}Ben Lonnqvist · CV · October 2026}
\fancyfoot[R]{\ttfamily\scriptsize\color{muted}\thepage/\pageref{LastPage}}

% Entry: year | body
\newlist{cventries}{itemize}{1}
\setlist[cventries]{leftmargin=21mm,labelwidth=21mm,labelsep=0pt,align=left,itemsep=0pt,topsep=0pt,parsep=0pt}
\newcommand{\entry}[2]{\begin{cventries}\item[{\ttfamily\small\color{muted}#1}] #2\end{cventries}\vspace{4pt}}
\newcommand{\role}[3]{\textbf{#1}\hfill{\small\color{muted}#2}\par{\small #3}\par}
\newcommand{\me}{\textbf{Lonnqvist, B.}}
\newcommand{\venue}[1]{\textit{#1}}
\newcommand{\lnk}[2]{\href{#1}{\ttfamily\footnotesize\textcolor{signal}{#2}}}
\setlist[itemize]{leftmargin=12pt,labelsep=5pt,itemsep=1pt,topsep=2pt,label={\textcolor{muted}{–}}}

\begin{document}\raggedright\sloppy

% ----- Header -----
{\cond\fontsize{30}{30}\selectfont BEN LONNQVIST}\par\vspace{5pt}
{\large Head of NeuroAI, \href{https://readwritelabs.com}{Read/Write Labs}}\par\vspace{6pt}
{\ttfamily\scriptsize\color{muted}
benhlonnqvist [at] gmail [dot] com\enspace·\enspace
\href{https://benlonnqvist.github.io}{benlonnqvist.github.io}\enspace·\enspace
\href{https://scholar.google.com/citations?user=be0tNVwAAAAJ}{Google Scholar}\enspace·\enspace
\href{https://github.com/benlonnqvist}{github.com/benlonnqvist}\enspace·\enspace
\href{https://x.com/lonnqvistben}{@lonnqvistben}}
\par\vspace{8pt}
{\color{ink}\rule{\linewidth}{1.2pt}}\par\vspace{7pt}
{\small Technical lead in NeuroAI. At Read/Write Labs I lead the NeuroAI team, and I architect the company's generative neuromodulation platform.
Research background in large-scale evaluation of AI against human vision and the brain, with multiple first-author papers
at top machine-learning venues.}

% ----- Experience -----
\cvsection{Experience}
\entry{2026–now}{\role{Head of NeuroAI}{}{Read/Write Labs (formerly Dandelion Science)}
\begin{itemize}\small
  \item Lead NeuroAI: set the technical direction for how brain models are built, validated and used to design targeted therapies for brain disorders.
  \item Architect and lead the development of the company's generative neuromodulation platform.
\end{itemize}}
\entry{2025–2026}{\role{AI Researcher}{}{Read/Write Labs (formerly Dandelion Science)}
\begin{itemize}\small
  \item Built and evaluated maintainable and sustainable computational brain models for the company's therapy-design platform.
  \item Led technical platform development.
\end{itemize}}
\entry{2020–2025}{\role{Doctoral Researcher}{Lausanne, Switzerland}{EPFL · NeuroAI Lab and Laboratory of Psychophysics}
\begin{itemize}\small
  \item Led a large-scale comparison of humans and 1,000+ AI models on contour integration (ICML 2025).
  \item Core contributor to the open-source Brain-Score platform (\href{https://github.com/brain-score}{github.com/brain-score}).
\end{itemize}}
\entry{2018–2020}{\role{Research Assistant \& Intern}{Aberdeen, UK}{School of Psychology, University of Aberdeen}
\begin{itemize}\small
  \item Deep-network models of visual crowding (first-author paper in \venue{Neural Networks}); Bayes-optimal eye-movement modelling and eye-tracking studies.
\end{itemize}}

% ----- Education -----
\cvsection{Education}
\entry{2020–2025}{\role{Ph.D., Computational Neuroscience}{EPFL, Lausanne}{Thesis: \textit{Differences in Visual Perception in Humans and Deep Neural Networks}\\
Nominated for the Best Thesis Prize\\
Advisors: Prof.\ Michael H.\ Herzog, Prof.\ Martin Schrimpf · Computer Science and Life Sciences departments}}
\entry{2016–2020}{\role{M.A.\ (Hons), Economics and Finance, First Class}{University of Aberdeen}{Thesis: \textit{Optimal search in multi-armed bandit problems}}}

% ----- Publications -----
\cvsection[* equal contribution]{Publications}
\cvsub{Conference papers}
\entry{2026}{Mehrer, J., \me, Mitola, A., Papale, P., Schrimpf, M. Model-guided microstimulation steers primate visual behavior.
\venue{International Conference on Learning Representations (ICLR)}. \lnk{https://arxiv.org/abs/2510.03684}{arXiv:2510.03684}}
\entry{2025}{\me, Scialom, E., Gokce, A., Merchant, Z., Herzog, M.\,H., Schrimpf, M. Contour integration underlies human-like vision.
\venue{International Conference on Machine Learning (ICML)}, PMLR 267:40290--40311. \lnk{https://proceedings.mlr.press/v267/lonnqvist25a.html}{proceedings.mlr.press}}
\entry{2024}{\me*, Wu, Z.*, Herzog, M.\,H. Latent noise segmentation: How neural noise leads to the emergence of segmentation and grouping.
\venue{International Conference on Machine Learning (ICML)}, PMLR 235. \lnk{https://proceedings.mlr.press/v235/lonnqvist24a.html}{proceedings.mlr.press}}
\cvsub{Journal articles}
\entry{2026}{Scialom, E., \me, Bornet, A., Herzog, M.\,H. Diagnostic features predicted by deep learning improve human object recognition in simulated prosthetic vision.
\venue{Journal of Neural Engineering}. \lnk{https://doi.org/10.1088/1741-2552/ae9227}{doi:10.1088/1741-2552/ae9227}}
\entry{2021}{\me, Bornet, A., Doerig, A., Herzog, M.\,H. A comparative biology approach to DNN modeling of vision: A focus on differences, not similarities.
\venue{Journal of Vision}, 21(10):17. \lnk{https://doi.org/10.1167/jov.21.10.17}{doi:10.1167/jov.21.10.17}}
\entry{2020}{\me, Clarke, A.\,D.\,F., Chakravarthi, R. Crowding in humans is unlike that in convolutional neural networks.
\venue{Neural Networks}. \lnk{https://doi.org/10.1016/j.neunet.2020.03.021}{doi:10.1016/j.neunet.2020.03.021}}
\cvsub{Preprints}
\entry{2022}{\me, Machiraju, H., Herzog, M.\,H. A comment on Guo et al.\ [arXiv:2206.11228]. \venue{arXiv}. \lnk{https://arxiv.org/abs/2208.01456}{arXiv:2208.01456}}
\entry{2020}{\me, Elsner, M., Hunt, A.\,R., Clarke, A.\,D.\,F. Modeling individual variation in visual search with reinforcement learning.
\venue{PsyArXiv}. \lnk{https://doi.org/10.31234/osf.io/suj28}{doi:10.31234/osf.io/suj28}}

% ----- Talks -----
\cvsection{Talks \& presentations}
\cvsub{Invited talks}
\entry{2019–2024}{Scene Grammar Lab, Goethe Universität Frankfurt (2024); INCC, Université Paris Cité \& CNRS Vision Group Symposium (2023); Psychology \& Computer Science, University of Essex (2019).}
\cvsub{Conference orals}
\entry{2019–2024}{Vision Sciences Society: perceptual grouping with latent noise (2024, main meeting and MODVIS); segmentation in noisy VAEs (2023, MODVIS); global information processing in feedforward networks (2022). Scottish Vision Group: object recognition in CNNs vs.\ humans (2019).}
\cvsub{Selected posters}
\entry{2021–2025}{ICML 2025 (contour integration); CCN 2024 (integration across object discontinuities); ECVP 2024 (grouping failures in DNNs, joint first author); VSS 2024 (categorical object properties, with Scialom et al.); VSS 2023 (segmentation in noisy VAEs); VSS 2021 (crowding in feedforward CNNs).}

% ----- Recognition & service (condensed) -----
\cvsection{Recognition \& service}
\entry{Press}{EPFL News front page (2025) and \textit{20 Minutes}, Switzerland (2025) on contour integration; \textit{Tekniikka\&Talous}, Finland (2025), interview on Latent Noise Segmentation.}
\entry{Awards}{Nominated for the Best Thesis Prize, EPFL; Rank Prize Foundation Optoelectronics Studentship (2019); M.A.\ with First-Class Honours (2020).}
\entry{Reviewing}{NeurIPS; CCN; ICML Workshop on LLMs and Cognition; NeurIPS Workshop on Behavioral ML. Journals: \textit{Neural Networks}, \textit{Neural Computation}, \textit{PLOS One}, \textit{Vision Research}.}
\entry{Teaching}{Teaching assistant, Understanding Statistics (M.Sc.) and Real Analysis (B.Sc.), EPFL 2020--2024. Supervised M.Sc.\ theses of M.\ van Holk, M.\ Ben Salah, C.\ Lévy-Fidel and Z.\ Wu, plus eight semester, summer and course projects.}

\end{document}
